\newcommand{\AbstractResults}{ApiaViz completed three of four routes, whereas
Sobel filtering with colour completed all four. ApiaViz had lower average
tracking error (11.5 versus 18.5\,cm, including failures). These results show
that deterministic spiking representations can support navigation, but do not
establish an overall advantage for ApiaViz over simple filtering.}

\newcommand{\ResultsParagraph}{ApiaViz reached the nest on three routes;
Sobel plus colour reached it on all four. Grayscale also completed three routes,
while raw colour, simple opponency and Ardin-style input each completed two.
ApiaViz's mean route
deviation was 11.5\,cm, compared with 18.5\,cm for Sobel and 33.1\,cm for grayscale.
Ardin-style input had a mean deviation of 66.4\,cm.
ApiaViz and Sobel produced similar activity levels: 4.95\% of KCs for
ApiaViz and 4.99\% for Sobel. Their different outcomes therefore occurred at
similar average sparsity.}

\newcommand{\DiscussionOpening}{The comparison separates two useful outcomes:
staying close to the taught route and eventually reaching the nest. ApiaViz
tracked three routes closely but failed on Ant~8. Sobel completed that route
and the other three, although its Ant~11 trial took 183 steps and had a mean
deviation of 53.6\,cm. ApiaViz may help tracking precision, but this experiment
does not show that it is a more reliable navigation input than Sobel filtering.}

\newcommand{\ConclusionParagraph}{A deterministic spiking navigation system
is feasible in this setting. ApiaViz reduced tracking error, while Sobel
produced the most consistent completion. Recovery after drift is therefore
the immediate priority for ApiaViz, followed by tests across more environments
and neural wiring seeds with the same matched controls.}
